% Courbes de titrage pH-métrique : méthode des tangentes et méthode de la dérivée (L'essentiel, chapitre 3)
\begin{tikzpicture}
\begin{axis}[otMesures,otAxesBout,name=tg,width=78mm,height=58mm,xmin=0,xmax=40,ymin=0,ymax=13,
  xtick={0,10,20,30,40},ytick={0,2,...,12},xlabel={$V_\mathrm{B}$ (mL)},ylabel={pH}]
\addplot[otNavy,line width=1pt] table[x index=0,y index=1] {figures/donnees/03-ethanoique-ph.dat};
\addplot[otMuted,line width=0.6pt,domain=8:23] {0.135*x+3.205};
\addplot[otMuted,line width=0.6pt,domain=16:32] {0.135*x+7.739};
\addplot[otGreenDark,line width=0.7pt,dashed,domain=11:29] {0.135*x+5.472};
\addplot[only marks,mark=*,mark size=2pt,otGreenDark] coordinates {(20.0,8.17)};
\draw[otGreenDark,line width=0.6pt,dash pattern=on 2pt off 1.5pt] (axis cs:20,8.17) -- (axis cs:20,0);
\draw[otGreenDark,line width=0.6pt,dash pattern=on 2pt off 1.5pt] (axis cs:20,8.17) -- (axis cs:0,8.17);
\node[pcLegende,anchor=north west,text=otGreenDark] at (axis cs:20.4,1.6) {$V_\mathrm{E}=20{,}0$ mL};
\node[pcLegende,anchor=south west,text=otGreenDark] at (axis cs:0.5,8.3) {pH$_\mathrm{E}\approx8{,}2$};
\node[pcLegende,anchor=west] at (axis cs:21.5,6.6) {E};
\node[pcLegendeG,anchor=north west,text=otNavy,align=left] at (axis cs:0.8,12.9) {méthode des\\tangentes parallèles};
\end{axis}
\begin{axis}[otMesures,otAxesBout,at={(tg.east)},anchor=west,xshift=17mm,width=78mm,height=58mm,xmin=0,xmax=40,ymin=0,ymax=7,
  xtick={0,10,20,30,40},ytick={0,2,...,6},xlabel={$V_\mathrm{B}$ (mL)},ylabel={$\mathrm{dpH}/\mathrm{d}V_\mathrm{B}$ (mL$^{-1}$)}]
\addplot[otNavy,line width=1pt] table[x index=0,y index=2] {figures/donnees/03-ethanoique-ph.dat};
\draw[otGreenDark,line width=0.6pt,dash pattern=on 2pt off 1.5pt] (axis cs:20,0) -- (axis cs:20,6.27);
\node[pcLegende,anchor=west,text=otGreenDark] at (axis cs:21.5,4.6) {maximum pour\\$V_\mathrm{E}=20{,}0$ mL};
\node[pcLegendeG,anchor=north west,text=otNavy] at (axis cs:0.8,6.9) {méthode de la dérivée};
\end{axis}
\end{tikzpicture}
\par\smallskip\legende{Titrage de 20,0 mL d'acide éthanoïque à \qty{1,00e-2}{mol.L^{-1}} par la soude à \qty{1,00e-2}{mol.L^{-1}}. Le saut de pH signale l'équivalence.}
